\newpage \begin{bclogo}[barre=none,logo={}]{Dissertation sur le sujet d'économie}
	\begin{center}
		\subsubsection*{Quelle politique structurelle pour un développement auto-entetenu des pays ouest-africains six décenies après les indépendances ?}\addcontentsline{toc}{subsection}{20240513 - UAO-M1-sujet (\'Eco)}
		%date - ISE 20... Sujet (Eco-1)
		%date - ISE 20... (Sujet 3)
	\end{center}
\end{bclogo}
%$$\star \star \star$$


%%\noindent \textbf{Définition des termes du sujet :}
%
%\begin{enumerate}[label=\theenumi\arabic* -]
%	\item 
%\end{enumerate}

%\begin{center}
%	\textbf{\underline{Introduction}} 
%\end{center}
	\hspace*{.5cm}L’analyse de Jiang (2011) met l’accent, entre autres, sur des relations de convergence et parvient à des conclusions intéressantes. Le premier résultat est que les transformations structurelles régionales ont contribuées à la convergence des produits intérieurs bruts par tête en termes réels des régions chinoises au cours de la décennie 1990. Le deuxième résultat indique que l’ouverture des régions favorise les transformations structurelles régionales. Ce bilan met bien en évidence, d'une part, le rôle majeur des transformations structurelles dans le développement de la Chine et dans la réduction des inégalités de revenus entre ces régions. D'autre part, il évoque l'effet positif de l'intégration régional sur les transformations structurelles. Cet exemple nous invite ainsi à réfléchir sur la politique structurelle à mettre en place dans les pays ouest-africains en vue de favoriser un développement autoentretenu de ces \'Etats,  soixante années après leurs indépendances. 
	
	\hspace*{.5cm}En premier lieu, La \guillemetleft\ politique structurelle\ \guillemetright\ désigne l'ensemble des transformations structurelles de l’économie – réglementation ou déréglementation, politique industrielle et politique sociale – opérées par les pouvoirs publics pour, ou bien redresser l'économie, ou bien pour l’adapter aux conditions économiques actuelles en vue d’atteindre, à long terme, le développement. En deuxième lieu, le \guillemetleft\ développement autoentretenu\ \guillemetright\ serait l’ensemble des changements qualitatifs induits par une croissance économique endogène résultant des forces internes du système économique – capital humain (niveau d’éducation et de santé de la population), du capital fixe (ensemble de biens de production de l’économie), du capital technologique (connaissances relatives à la production), et capital public (infrastructures financées par la puissance publique). En dernier, le \guillemetleft\ pays ouest-africains\ \guillemetright\ pourrait aussi désigner dans une certaine mesure les états membres de la commission économique des États de l’Afrique de l’ouest (CEDEAO) créée en 1975 à Lagos, Nigéria. 
	
	\hspace*{.5cm} Il subsisterait une politique structurelle capable d’engendrer des changements qualitatifs à long terme des économies ouest-africains. Plus précisement, il est encore possible pour les pouvoirs publics des \'pays de la CEDEAO, six décennies après leurs indépendances, d'effectuer des transformations structurelles de leurs économies afin d'entrainer leur développement autoentretenu. Ce qui implique, dès lors, l'existence d'un moyen d'action visant la réalisation de ces transformations. De prime abord, une réglementation qui favoriserait aux entreprises locales l'accès aux différents marchés semble servir de point de départ. ce faisant, en protégeant par exemple ces entreprises de l'influence de la mondialisation, celles-ci sont plus en voie de devenir des champions nationaux. De surcroit, une politique industrielle ayant comme socle l'innovation et la R\&D serait un élément incontournable. Car, l'innovation participe, selon Joseph Schumpeter, à la mutation industrielle nécessaire à la révolution des structures économiques. \mbox{\emph{Cf} \underline{Capilisme, socialisme et démocratie}}. Aussi, une politique sociale inclusive serait à prendre en compte puisqu'elle favorise à la fois la productivité d'une économie et le développement humain. Comment peut-on justifier l'existence actuelle d'une politique structrurelle à même d'engendrer un développement autoentretenu des \'Etats membres de la CEDEAO après six décennies sans qu'ils n'aient pu garantir leur développement ? Peut être n'ont-ils pas porté une attention particulière simultanément à leurs entreprises locales et au marché régional à fortes potentialités de la CEDEAO comme l'auraient fait la majorité de BRICS (Brésil, Russie, Inde, Chine et Afrique du Sud) et les \'Etats Européens. Faut-il en déduire que les modèles d'industrialisation adopté par ces \'Etats seraient aussi un exemple à suivre pour les \'Etats de l'Afrique de l'ouest ? Cette hypothèse semble crédible vu que leurs modèles leur ont fourni des résultats satisfaisants. Quid de la politique sociale à adopter alors ?
	
	\hspace*{.5cm} Quelle combinaison de transformations structurelles adoptée par les économies actuelles de la CEDEAO favoriserait leur développement auto-entretenu ?\\
	
	\hspace*{.5cm}Si la réglementation des économies actuelles de la CEDEAO en vue favoriser l’accès aux marchés locaux, régionaux et internationaux aux PME/PMI est nécessaire pour atteindre le développement autoentretenu de ces \'Etats (I), la mise en place d'une politique industrielle basée sur l'innovation et la R\&D ne serait pas à omettre (II), outre une politique d'intégration sociale (III). $$\star\star\star$$